\clearpage
\section{Parametric polymorphism and conservative erasure}
\label{sec:generics}

This section puts the source language of Section~\ref{sec:generic-source}
and the logic of Section~\ref{sec:generic-logic} to work. The checked
definitions and proofs are in
\path{example/generics/rocq/generic_examples.v}. The program defines
\texttt{genericId<T>}, \texttt{Pair<A, B>}, and \texttt{Box<T>} with
\texttt{get} and \texttt{map<U>}. The same definitions feed the generic
printer. Each contract below is uniform in its element predicates and
source type arguments; it is not a collection of integer-only proofs.

\subsection{From a declaration to the program being proved}

Identity is a particularly small example of a generic declaration:
\begin{lstlisting}
Definition identity_body := GEVar "x".
Definition identity_fn :=
  GFunction "genericId" None ["T"]
    [("x", GTVar "T")] (GTVar "T") identity_body.
\end{lstlisting}
The declaration has no owner, binds one type parameter and one value
parameter, and returns that value. Its source call is
\texttt{GECall "genericId" [t] [GEVal v]}. Linking produces a recursive
core value with an internal tuple parameter; its body projects argument
zero and returns it. The annotations \texttt{T} and \texttt{t} are absent
from that value. The printer instead retains them, producing
\begin{lstlisting}
public func genericId<T>(x: T): T {
    return x
}
\end{lstlisting}
Thus there are two uses of one structured source: erasure supplies the
program for the Iris proof, and printing supplies the native generic
declaration. The mathematical notation \(\GWpre\) below fixes
\texttt{generic\_layouts} and \texttt{generic\_runtime}; source-like
calls inside it abbreviate their explicit \texttt{gexpr} constructors.

\subsection{Identity: one proof for every representation}

For every predicate \(R:\mathsf{val}\to\mathsf{iProp}\,\Sigma\),
type argument \(t\), and value \(v\), \texttt{identity\_spec} establishes
\[
R(v)\;\vdash\;
\GWpre\,(\mathsf{genericId}\langle t\rangle(v))
  \,\{w.\;\pure{w=v}\sep R(w)\}.
\]
The result is the exact input value, and the client's representation
resource is returned. No persistence assumption on \(R\) is used.
The proof script is short enough to display in full:
\lstinputlisting[firstline=75,lastline=85]{example/generics/rocq/generic_examples.v}

\texttt{wp\_generic\_pack\_acc} evaluates the internal argument tuple.
\texttt{wp\_app} enters the linked function; \texttt{wp\_struct\_load}
extracts its parameter, and \texttt{wp\_let} substitutes it. The final
\texttt{iFrame} returns the untouched \(R(v)\). This explains the
uniformity: the body neither inspects the element nor requires any facts
about its representation. The theorem is also independent of whether
\(t\) is \texttt{Int64}, \texttt{Bool}, or a nested named application.
Well-formed source types and native typing are separate checks from this
semantic statement about values.

For example, one may take \(R(w)=\pure{w=42}\),
\(R(w)=\pure{w=\mathsf{true}}\), or a predicate owning an object's
fields. These are choices of logical predicates, not runtime type tests.
The closed theorem \texttt{generic\_main\_adequate} applies this proof
to the integer call returning 42, as described in
Section~\ref{sec:generic-adequacy}.

\subsection{Pairs: independent parameters and projection resources}

\texttt{Pair<A, B>} is a struct with ordered fields \texttt{first: A}
and \texttt{second: B}. Construction erases to a two-value tuple; no heap
allocation is required. Its representation predicate is
\[
\mathsf{PairRep}(R,S,p)\;\triangleq\;
\exists a,b.\;\pure{p=\langle a,b\rangle}\sep R(a)\sep S(b).
\]
The predicates \(R\) and \(S\) are independent, just like the type
parameters \(A\) and \(B\). The constructor theorem is
\path{pair_construct_spec}:
\[
R(a)\sep S(b)\;\vdash\;
\GWpre\,(\mathsf{Pair}\langle t_a,t_b\rangle(a,b))
  \,\{p.\;\mathsf{PairRep}(R,S,p)\}.
\]
The proof only evaluates the fields and packages the two owned predicates:
\lstinputlisting[firstline=87,lastline=93]{example/generics/rocq/generic_examples.v}

The source functions \texttt{firstOf<A, B>} and \texttt{secondOf<A, B>}
project indexes zero and one of the explicitly named \texttt{Pair} layout.
Their contracts are
\[
\begin{aligned}
\mathsf{PairRep}(R,S,p)&\vdash
 \GWpre\,(\mathsf{firstOf}\langle t_a,t_b\rangle(p))\\
 &\qquad\{w.\;\exists b.\;\pure{p=\langle w,b\rangle}\sep R(w)\sep S(b)\},\\
\mathsf{PairRep}(R,S,p)&\vdash
 \GWpre\,(\mathsf{secondOf}\langle t_a,t_b\rangle(p))\\
 &\qquad\{w.\;\exists a.\;\pure{p=\langle a,w\rangle}\sep R(a)\sep S(w)\}.
\end{aligned}
\]
The equality identifies the returned component, and each component
resource occurs once. In particular the first contract does not return
\(\mathsf{PairRep}(R,S,p)\sep R(w)\), which would duplicate an arbitrary
resource. The checked proof of the first projection is
\lstinputlisting[firstline=97,lastline=109]{example/generics/rocq/generic_examples.v}
The second proof uses the same argument at index one. Copying the struct
value copies its stored values; it does not create another ownership
permission for an object referenced by either value.

\subsection{Boxes: allocation and reading}

\texttt{Box<T>} is a final class with one immutable source field,
\texttt{value: T}. Its runtime representation is an object reference.
Writing \(o\ptso[0]v\) for \texttt{ObjId o} owning field zero, define
\[
\mathsf{BoxRep}(R,p)\;\triangleq\;
\exists o,v.\;\pure{p=\mathsf{obj}(o)}\sep o\ptso[0]v\sep R(v).
\]
This owns both the box's field and whatever resources describe its
element. The field's source immutability does not make either ownership
predicate persistent in this Iris model.

\texttt{box\_construct\_spec} and \texttt{box\_get\_spec} state
\[
\begin{aligned}
R(v)&\vdash
 \GWpre\,(\mathsf{Box}\langle t\rangle(v))
   \,\{p.\;\mathsf{BoxRep}(R,p)\},\\
\mathsf{BoxRep}(R,p)&\vdash
 \GWpre\,(p.\mathsf{get}())\\
 &\qquad\{w.\;\exists o.\;\pure{p=\mathsf{obj}(o)}
      \sep o\ptso[0]w\sep R(w)\}.
\end{aligned}
\]
Construction uses \texttt{wp\_new} to obtain the fresh field ownership and
combines it with \(R(v)\). Reading uses \texttt{wp\_field\_load}, which
returns that same field ownership. Neither operation duplicates \(R\).
The implementation proofs are
\lstinputlisting[firstline=125,lastline=145]{example/generics/rocq/generic_examples.v}

\subsection{A type-changing method and its callback contract}
\label{sec:generic-map}

In \texttt{Box<T>.map<U>}, \texttt{T} is fixed by the receiver and
\texttt{U} is chosen at the call. The callback has type \texttt{(T) -> U},
and the result has type \texttt{Box<U>}. The method reads the old value,
calls the callback once, and allocates a new box for its result.
Its erased body, after parameter substitution, is
\[
 \mathsf{New}\;[]\;
   [\mathsf{App}\;(\mathsf{Val}\;f)
      \;(\mathsf{FieldLoad}\;(\mathsf{obj}(o))\;0)].
\]
For input and output predicates \(R,S\), the required callback contract is
\[
\mathsf{Callback}(R,S,f)\;\triangleq\;
 \forall a.\;R(a)\wand\Wpre\,(f\,a)\,\{S\}.
\]
It provides one invocation, consuming the input resource and supplying the
output resource. It does not assert that the callback is pure or that its
contract is duplicable. The additional Coq premise
\(\forall x,w.\;\mathsf{subst\_val}\;x\;w\;f=f\) requires a closed
callback value, so unpacking the method's parameters cannot capture free
variables in it.

Under that premise, \texttt{box\_map\_spec} is
\[
\begin{aligned}
o\ptso[0]v\sep R(v)\sep\mathsf{Callback}(R,S,f)
 &\vdash \GWpre\,
   (\mathsf{obj}(o).\mathsf{map}\langle u\rangle(f))\\
 &\qquad\{p.\;o\ptso[0]v\sep\mathsf{BoxRep}(S,p)\}.
\end{aligned}
\]
The old box keeps the same field value. Its element resource \(R(v)\)
has gone to the callback; the result box receives \(S\). Thus the
postcondition deliberately retains the old field, rather than asserting
a second \(\mathsf{BoxRep}(R,\mathsf{obj}(o))\).

The complete checked statement and proof make the ownership flow explicit:
\lstinputlisting[firstline=150,lastline=179,basicstyle=\ttfamily\footnotesize]{example/generics/rocq/generic_examples.v}
After tuple unpacking, \texttt{wp\_field\_load} obtains the element while
preserving \texttt{Ho}, the old field assertion. The subproof supplied
with \texttt{[HR Hf]} gives only the element resource and callback contract
to the callback. \texttt{Ho} is framed across it. Finally,
\texttt{wp\_new} creates \texttt{Ho'}, and the proof packages
\texttt{Ho'} with \texttt{HS} into the result representation.

For a concrete ownership example, let \(R=\mathsf{BoxRep}(P,-)\)
describe an inner box and let the callback return that inner box unchanged,
so \(S=R\). Mapping an outer box produces a new outer object containing
the same inner reference. The inner field ownership is transferred to the
new outer representation. The old outer field is still owned, but one
cannot claim that both outer boxes separately own the same inner box.
If \(R\) is persistent, such as a pure integer-value assertion, the usual
Iris persistence rules allow it to be retained. That is an extra property
of the chosen predicate, not an assumption hidden in \texttt{map}.

\subsection{Concrete and nested instances}

The runtime suite uses both primitive and nested types. The following
fragments illustrate the calls made possible by the declarations:
\begin{lstlisting}
let n: Int64 = genericId<Int64>(42)
let b: Bool = genericId<Bool>(true)
let p = Pair<Int64, Bool>(8, true)
let first: Int64 = firstOf<Int64, Bool>(p)
let second: Bool = secondOf<Int64, Bool>(p)

let source = Box<Int64>(5)
let target: Box<Bool> = source.map<Bool>(
    { x: Int64 => 0 <= x })

let nested = genericId<Box<Pair<Int64, Bool>>>(
    Box<Pair<Int64, Bool>>(p))
let result: Box<Box<Bool>> = nested.map<Box<Bool>>(
    { x: Pair<Int64, Bool> => Box<Bool>(x.second) })
\end{lstlisting}
Here the method parameter is instantiated to \texttt{Bool} in the first
map and to \texttt{Box<Bool>} in the second; it differs from the receiver's
parameter in both cases.

One can express the nested representation by setting
\(R_8(v)=\pure{v=8}\), \(R_{\mathsf{true}}(v)=\pure{v=\mathsf{true}}\),
and \(P(p)=\mathsf{PairRep}(R_8,R_{\mathsf{true}},p)\).
Then \(\mathsf{BoxRep}(P,b)\) describes the nested input box.
For the last callback, take
\(S(w)=\mathsf{BoxRep}(R_{\mathsf{true}},w)\): projecting the second
component and allocating its Boolean box supplies the callback obligation.
The map rule returns the outer field of the old box and
\(\mathsf{BoxRep}(S,\mathit{result})\). This is how the uniform
contracts can be composed; the particular native test is a runtime check,
not a separately named closed adequacy theorem for this whole fragment.

\texttt{repeatIdentity<T>} also exercises direct generic recursion.
Its native tests instantiate it at integers, booleans, and pairs.
The suite checks recursion and printing; this example does not supply a
separate general termination proof for that function.

\subsection{Generated source and extraction checks}
\label{sec:generic-extraction}

The emitted program below is read directly from the generated file.
The small identity-lambda wrappers in field accesses force the native
compiler to check the annotated receiver type. Similar wrappers check
lambda result annotations. They are not used to select layouts by value:
the source explicitly names the declaration and field index.
\lstinputlisting{example/generics/extracted/generics.cj}

The checked front end returns \texttt{Extracted} only after structural
validation; otherwise \texttt{Rejected} contains diagnostics. It checks
declaration names, type-variable scopes, type and value arities, named
receiver layouts, field indexes, and supported forms. For example, the
regression suite includes these failures:
\begin{lstlisting}
genericId with no type argument:
  wrong type arity: genericId
missing<Int64>(1), with no such declaration:
  unresolved function (declare dependencies first): missing
Box<Bool>() with no field argument:
  wrong constructor arity: Box
Box<T>.map<T> rebinding its owner's parameter:
  invalid, shadowed, or duplicate type parameters
\end{lstlisting}
This listing shows diagnostics rather than runnable Cangjie source.
Checking type compatibility of ordinary arguments and results is left
to Cangjie compilation. Tests deliberately compile invalid receiver and
lambda annotations and require rejection. Distinct declarations with
identical bodies or equal-sized layouts have separate explicit references;
the printer does not choose between them by inspecting those bodies or sizes.

Evaluation order is part of this boundary. Calls and constructors keep
left-to-right argument order; the receiver precedes method arguments.
CJR's Boolean conjunction evaluates both operands, so the generic printer
uses a helper call to preserve that evaluation before applying native
\texttt{\&\&}. The \texttt{strictAnd} runtime case checks the right-hand
callback's effect even when the left-hand value is false.

\subsection{Validation and the scope of the guarantees}

\path{example/generics/rocq/generic_tests.v} checks simultaneous and nested
type substitution, type-parameter scopes, instantiation erasure, local
substitution under binders, and extraction rejection cases in the Rocq
kernel. \path{example/generics/test/generics_test.cj} supplies ten native
test cases, including heterogeneous pairs, type-changing maps, nested
instances, identical-body declarations, and direct recursion. Five
compiler-rejection tests in \path{tests/test_generic_extraction.py}
check unrelated receiver layouts, incorrect receiver type arguments, and
incorrect lambda result types, buffer receiver annotations, and buffer
write types against the Cangjie 1.0.5 toolchain. Section~\ref{sec:poly-arrays}
extends the examples with managed polymorphic arrays and a paired benchmark.

\texttt{make verify} includes the erasure laws, generic logic,
representation contracts, and generic adequacy results in the assumption
audit. \texttt{make test} includes the native generic and existing runtime
suites. Existing monomorphic examples retain their original printer and
proofs; the generic addition does not migrate their data structures.

The result is uniform Iris specifications for the supplied implementations
and safe execution of erased programs under their proved preconditions.
There is no general typing-soundness or relational-parametricity theorem.
The printer and native compiler remain in the translation trust boundary,
and the unbounded-integer model still requires care at \texttt{Int64}
boundaries. Interface bounds, inheritance, higher-rank polymorphism, and
generic raw-memory allocation are outside the implemented source fragment.
